\section{Illustrative dialogues for the ESI} \label{app:cases}

This appendix illustrates how investor tone, managerial tone, and reply substantiveness enter the interpretation of a dialogue. A single example cannot reproduce the daily ESI for a firm because the index also uses question attention and aggregates all interactions for that firm on the same day. The sentiment scores below range from $-1$ to $1$ and are manual assessments for illustration. They are not model scores. The interpretations describe the direction of each tone signal and the substantiveness adjustment. They do not assign an exact attention weight or investor state weight.

\subsection{Scenario A: Negative investor tone and positive managerial tone}

The first two examples begin with a strongly negative investor assessment and receive a positive managerial response. The replies differ in substantiveness. ESI assigns more weight to managerial tone when the reply supports its assessment with specific information.

\noindent \textbf{D.1. Example 1} \\
\textbf{巴比食品} [Stock Code: 605338, 2023]. \textit{Illustrative scores: Investor sentiment is $-0.90$. Managerial sentiment is $+0.35$. Reply substantiveness is high. The positive managerial signal receives little discount from the substantiveness adjustment.} \\
\textbf{Question:} Why did the company's share price fall sharply? Has the company encountered a serious problem? First quarter net profit increased by 27 times. What caused the decline in the share price? (公司股价为什么暴跌，是暴雷了吗，一季度净利润不是增长27倍，到底是出了什么事？) \\
\textbf{Reply:} The company states that its operations are sound and that its store and group meal businesses continue to grow. First quarter net profit was about RMB 41 million, compared with RMB 1.45 million in the same period of 2022. The increase mainly reflects a low comparison base and changes in the share price of Dongpeng Beverage. Excluding nonrecurring items, first quarter net profit was about RMB 20.15 million, down 48\% from RMB 38.80 million in the same period of 2022. The company attributes this decline to slow revenue growth and higher depreciation, personnel, marketing, and equity incentive expenses. It expects revenue growth and net profit excluding nonrecurring items to improve as the temporary effects dissipate. (尊敬的投资者，您好。公司目前经营情况良好，门店业务和团餐业务保持正常增长。一季度实现净利润约4,100万元，同比2022年一季度145万元增长较多主要系去年同期基数偏低，同时公司净利润每季度受东鹏饮料股价变动影响。扣除此类非经常性损益外，公司一季度实现扣除非经常性损益后的净利润约2,015万元，同比2022年一季度3,880万元下跌48\%。公司一季度扣非净利润下跌主要原因是收入增速较慢，但为扩大业务规模而准备的包括南京工厂在内的资产折旧、人员成本、各地加速开店的营销费用支出，以及今年实施员工股权激励计划须确认的股份支付费用同比增加较多所致。一季度收入增速较慢，主要系公司主要收入来源的门店业务增长相对缓慢。原因是2023年1-2月份春节假期前后，门店歇业周期和门店歇业数量相较往年增多，门店实际营业天数不足，拖累了公司整体业绩。随着不利因素消除，预计公司收入增速将进一步增长，相应的人员成本、资产折旧与一季度相比会保持金额稳定，公司扣非后净利润将会显著改善。谢谢！)
\vspace{0.4cm}

\noindent \textbf{D.2. Example 2} \\
\textbf{利德曼} [Stock Code: 300289, 2024]. \textit{Illustrative scores: Investor sentiment is $-0.70$. Managerial sentiment is $+0.30$. Reply substantiveness is low. The positive managerial signal receives a large discount from the substantiveness adjustment.} \\
\textbf{Question:} Investors remain concerned about the risk of a special treatment designation because the company's performance has remained weak for several years and has alternated between small profits and losses. What measures will management take to reverse this condition? (董秘你好，近期不断有投资者担心公司被st而退市，数次对公司提出质疑，董秘都以公司治理规范，不存在st退市风险，但是投资者看到的是公司近几年业绩持续低迷，虽然公司年度业绩在公司管理层努力下微利与亏损交替发生，暂时不存在连续二年亏损而st退市，但是公司业绩持续低迷是不争的事实，请问公司今后采取什么措施扭转不利局面？) \\
\textbf{Reply:} The company will maintain prudent operations, improve operating and managerial efficiency through cost reductions, and explore new businesses and sources of profit. (您好！公司秉持稳健经营理念，通过降本增效等系列举措，提升运营效率和管理效能，并积极探索新的业务和利润增长点。感谢您的关注！)

\subsection{Scenario B: Positive investor tone and positive managerial tone}

The third example begins with a favorable investor assessment. Management confirms this assessment with financial results and an operating explanation. The managerial contribution remains positive, but it receives less state weight because the investor already expresses an optimistic view.

\noindent \textbf{D.3. Example 3} \\
\textbf{中国船舶} [Stock Code: 600150, 2024]. \textit{Illustrative scores: Investor sentiment is $+0.75$. Managerial sentiment is $+0.70$. Reply substantiveness is high. The managerial contribution to ESI is positive and receives less state weight than an equally positive reply following a negative question.} \\
\textbf{Question:} Congratulations on the company's strong performance in the first half of the year. Which business segments or projects drove this growth? Can these sources of growth continue? (恭喜公司在上半年取得良好业绩。请问这些显著增长的背后，主要有哪些具体的业务板块或项目驱动？未来这些增长动能是否可持续？) \\
\textbf{Reply:} First half revenue was RMB 36.017 billion, an increase of 17.99\%. The company attributes this growth mainly to an increase in completed vessel deliveries. Revenue from shipbuilding, repair, and marine engineering was RMB 34.446 billion, an increase of 22.39\%. Management also cites a shift toward vessels with higher value added and stronger orders during the recovery of the shipping and shipbuilding industries. (感谢您的关注。2024年上半年，公司践行“三赢”理念，聚焦主责主业，强化价值创造，着力提升运行质量，加快培育船舶新质生产力，不断增强核心能力、提升核心竞争力，实现营业收入360.17亿元，同比增长17.99\%，主要系完工交船较上年同期大幅增加所致。其中：船舶造修及海洋工程业务营业收入344.46亿元，同比增长22.39\%。公司作为国内规模最大、技术最先进、产品结构最全的造船旗舰上市公司，在研发、人才、规模、产品结构等方面有明显优势，持续突破高附加值船舶，船型结构持续优化，在本轮航运和造船业复苏的背景下，取得了较好的接单成果，打开了长期可持续发展的空间。)

\subsection{Scenario C: Negative investor tone and negative managerial tone}

The final example shows how a substantive reply can transmit adverse information. Substantiveness increases the weight assigned to the managerial signal. The negative tone of the reply determines the direction of its contribution.

\noindent \textbf{D.4. Example 4} \\
\textbf{ST工智} [Stock Code: 000584, 2023]. \textit{Illustrative scores: Investor sentiment is $-0.80$. Managerial sentiment is $-0.85$. Reply substantiveness is high. The negative managerial signal receives little discount from the substantiveness adjustment.} \\
\textbf{Question:} What caused the company's losses over the past two years? Does management expect another loss this year? Is the company at risk of delisting? (请问贵公司近两年亏损的原因是什么？今年还会继续亏损吗？会退市吗？) \\
\textbf{Reply:} The company reports that its largest business, automotive manufacturing, has faced semiconductor shortages, high raw material prices, cash flow pressure, weak downstream demand, delayed or cancelled production line projects, and intense competition. These conditions reduced new projects and project margins. Financial distress among several automaker customers also increased receivable impairments. Based on its 2022 earnings forecast, the company expects net profit for 2022 to be negative. (尊敬的投资者，您好！近年来公司业务占比最大的汽车制造业面临了芯片短缺、原材料价格持续高位、现金流等不利因素影响。自2018年四季度下游汽车行业低迷，行业出现汽车整车厂新生产线投资项目暂缓或停止的情况，新增产线项目大幅减少，行业竞争激烈，导致公司签订的项目毛利率大幅下降公司业绩受到影响，尤其是2021年下半年包括南京知行汽车、宝沃汽车、华晨汽车、宝能汽车等在内的部分国内汽车整车制造厂商或进入被申请破产流程或申请破产或出现重大经营风险，导致公司报告期内的应收款项减值大幅增加。根据《江苏哈工智能机器人股份有限公司2022年度业绩预告》，预计公司2022年度净利润为负值。感谢您对公司的支持!)
